\section{Математический аппарат поиска путей с временной зависимостью}

Использование динамических функций сопротивления сети превращает классическую задачу маршрутизации на статическом графе в гораздо более сложную задачу поиска кратчайшего пути с временной зависимостью. Вес ребра перестает быть константой и становится функцией от абсолютного времени прибытия агента на данное ребро.

Динамический поиск кратчайших путей и кольцевой буфер прогнозирования.
Для моделирования динамики заторов в будущем горизонт прогнозирования транспортной обстановки (2 часа) разбивается на 24 дискретных временных интервала по 5 минут (300 секунд) каждый. 

Во время обхода графа алгоритмом динамического поиска, стоимость проезда по ребру рассчитывается не на основе текущего глобального времени симуляции, а на основе ожидаемого времени прибытия агента на начало данного ребра. 

Выбор конкретного временного интервала (индекса массива, хранящего прогнозируемый объем трафика) осуществляется через безблокировочный механизм кольцевого буфера. Индекс интервала t\_idx рассчитывается с использованием целочисленной арифметики и оператора деления по модулю:
\begin{equation}
    t\_idx = \left( \left\lfloor \frac{arrival\_time}{300} \right\rfloor + head\_index \right) \bmod 24
\end{equation}
где arrival\_time — абсолютное симулируемое время прибытия агента на сегмент, а head\_index — указатель на начало кольцевого буфера. Асинхронный менеджер времени периодически (каждые 5 минут) смещает указатель head\_index на единицу вперед, одновременно обнуляя самую старую ячейку в хвосте буфера через высокоскоростную инструкцию memset \cite{drepper2007memory}. Это позволяет роутеру беспрепятственно читать прогнозируемые объемы трафика без захвата мьютексов \cite{butorin2026thread}.

Доказательство допустимости эвристики на основе опорных вершин для динамических графов.
Классический алгоритм Дейкстры гарантированно находит кратчайший путь, но обходит слишком много нерелевантных узлов, распространяясь во все стороны графа концентрическими кругами \cite{cherkassky1994shortest}. Для ускорения поиска применяется алгоритм направленного поиска, который сокращает область поиска в сторону цели с помощью эвристической функции H(u, v) \cite{goldberg2004alt}. 

Чтобы алгоритм гарантированно находил математически оптимальный путь, эвристика должна быть допустимой, то есть она никогда не должна переоценивать истинное расстояние до цели \cite{goldberg2004alt}. В статичных сетях часто используется прямолинейное расстояние, однако в городских лабиринтах оно крайне неэффективно \cite{fuchs2010preprocessing}. В разработанной системе применена эвристика на основе опорных вершин (маяков) и неравенства треугольника.

Для набора маяков L и статических весов свободного потока T\_{free} эвристическое расстояние от текущего узла u до цели v вычисляется как максимальная разность расстояний до маяка \cite{goldberg2004alt}:
\begin{equation}
    H_{alt}(u, v) = \max_{L} |T_{free}(L, v) - T_{free}(L, u)|
\end{equation}
В силу неравенства треугольника, эта разность расстояний никогда не превышает прямое расстояние между u и v на пустом графе: 
\begin{equation}
    H_{alt}(u, v) \le T_{free}(u, v)
\end{equation}

Ключевой теоретический вопрос заключается в том, сохраняет ли данная эвристика свою допустимость при возникновении заторов, когда веса ребер возрастают.
Так как динамический штраф за затор всегда строго неотрицателен ($\Delta W_{dynamic} \ge 0$), реальное время проезда по нагруженной сети $W_{real}(u, v, t)$ всегда больше или равно времени проезда по пустой сети $T_{free}(u, v)$: 
\begin{equation}
    W_{real}(u, v, t) \ge T_{free}(u, v)
\end{equation}

Объединяя оба неравенства, получаем строгое математическое доказательство допустимости эвристики в любой момент времени $t$:
\begin{equation}
    H_{alt}(u, v) \le T_{free}(u, v) \le W_{real}(u, v, t)
\end{equation}

Данное доказательство имеет фундаментальное значение: эвристика, рассчитанная только один раз на этапе предобработки пустого графа, абсолютно легитимна и применима в условиях сильных заторов. Она никогда не переоценит реальную стоимость пути, гарантируя математическую точность работы роутера без необходимости инвалидации и пересчета матрицы маяков.

Однако при экстремальном росте заторов и взрывном увеличении весов ребер возникает иной негативный эффект~--- деградация эффективности поиска вследствие снижения точности эвристической оценки. Поскольку эвристика $H_{alt}(u, v)$ опирается на статические свободные веса $T_{free}$, при многократном росте динамических весов ребер под BPR-нагрузкой реальное накопленное расстояние $g(u)$ начинает колоссально превосходить эвристическую оценку: $g(u) \gg H_{alt}(u, v)$. В целевой функции выбора вершины:
\begin{equation}
    F(u) = g(u) + H_{alt}(u, v)
\end{equation}
слагаемое $H_{alt}(u, v)$ превращается в математический шум. Эвристический алгоритм $A^*$ фактически деградирует до классического алгоритма Дейкстры, теряя направляющие свойства и начиная раскрывать избыточные вершины кругами (до 57~тыс. вершин на коротких запросах), что катастрофически снижает производительность ядра (QPS).

Для преодоления этой проблемы на уровне C++ ядра в классе \texttt{CostCalculator} реализован \textbf{механизм ограничения деградации направляющих свойств эвристической функции (индекса эффективности поиска) в условиях экстремального динамического изменения топологии весов}. Данный механизм осуществляет жесткое ограничение динамического BPR-штрафа на уровне 10-кратного базового времени проезда свободного потока:
\begin{equation}
    W_{total}(u, v, t) = T_{free}(u, v) + \min\left( \Delta W_{dynamic}(u, v, t),\, 10 \cdot T_{free}(u, v) \right)
\end{equation}
Такое ограничение удерживает разрыв между реальной и статической стоимостью путей в пределах одного десятичного порядка, сохраняя направляющие эллиптические фронты поиска алгоритма ALT и не допуская деградации $A^*$ до Дейкстры в перегруженных зонах, при этом гарантируя физическую адекватность маршрутизации.


Ограничения двунаправленного поиска в динамических транспортных сетях.
Интеграция механизмов двунаправленного эвристического поиска ALT в разработанное вычислительное ядро мезоскопического симулятора принципиально нереализуема в условиях динамического пошагового бронирования емкости.

В классических статических графах двунаправленный поиск Дейкстры или $A^*$ удваивает скорость за счет одновременного распространения двух волн: прямой от источника $s$ и обратной от цели $t$. Однако при переходе к мезоскопической модели с предиктивными корзинками объемов, где веса ребер зависят от времени и бронируются агентами наперед, обратный поиск полностью ломается вследствие двух факторов:
\begin{enumerate}
    \item \textbf{Отсутствие временного якоря на целевой вершине $t$.} Прямой поиск стартует в жестко известный момент времени симуляции $T_{start}$ и, двигаясь по стреле времени вперед, вычисляет точный кумулятивный момент времени прибытия (ETA) на каждое последующее ребро. Обратная волна должна распространяться из вершины $t$ в прошлое. Однако точное время прибытия агента в точку $t$ остается неизвестным до момента полного завершения расчета кратчайшего пути. Инициализация обратного поиска фиктивным временем приводит к считыванию заторов из случайных временных интервалов буфера, что полностью инвалидирует расчеты.
    \item \textbf{Динамический характер бронирования.} В отличие от классических задач маршрутизации с временной зависимостью (TDSP), где профили заторов статичны, в разработанной системе агенты активно модифицируют состояние графа непосредственно в процессе расчета маршрута, инкрементируя объемы в корзинках. Обратная волна физически не способна осуществить корректное резервирование емкости, так как не обладает информацией о плотности потока, формируемой агентами, движущимися в прямом направлении.
\end{enumerate}

Таким образом, алгоритм обратного распространения волны, разворачивающийся от целевой вершины к источнику, в условиях динамического квантования времени лишен темпорального детерминизма. Это теоретически обосновывает выбор однонаправленного алгоритма ALT с высокопроизводительной векторной SoA-очередью приоритетов в качестве оптимального решения для разработанного вычислительного ядра.


Принцип сохранения порядка очередности движения и линейная интерполяция времени.
Дискретизация времени на 5-минутные слоты приводит к тому, что функция стоимости ребра от времени становится ступенчатой. Это грубо нарушает важнейший физический принцип маршрутизации — свойство сохранения порядка очередности движения (FIFO) \cite{goldberg2005external}. Нарушение этого свойства означает парадоксальную ситуацию: автомобиль, выехавший позже, может обогнать выехавший ранее автомобиль только потому, что его время прибытия попало в соседний временной интервал с резко меньшим штрафом.

Математическим условием соблюдения принципа очередности является непрерывность функции стоимости $W_{total}(t)$ и ограничение скорости убывания веса ребра: она не должна превышать скорость течения самого времени (условие Липшица): 
\begin{equation}
    \frac{\partial W_{total}(t)}{\partial t} \ge -1
\end{equation}

Для устранения ступенчатости и строгого соблюдения принципа очередности в разработанной системе применен механизм линейной интерполяции времени на границах временных интервалов. Вместо интерполяции итоговых весов, для сохранения высокой производительности интерполируются непосредственно вычисленные штрафы из двух соседних интервалов P\_1 (текущий интервал) и P\_2 (следующий интервал):
\begin{equation}
    \Delta W_{dynamic}(t) = \left( P_1 + \frac{(P_2 - P_1) \cdot (arrival\_time \bmod 300)}{300} \right) \gg 20
\end{equation}

В сочетании с макроскопическим законом сохранения транспортных средств (который физически лимитирует максимальный перепад объемов трафика между соседними интервалами) и алгоритмом предпространственной нарезки графа, механизм линейной интерполяции математически гарантирует выполнение принципа очередности во всей моделируемой транспортной сети.
